While the Ising model is not a perfect representation of the quantum mechanics in ferromagnets, it is impressive that the model predicts a phase transition at the critical temperature, just as real permanent magnets will break down above their Curie temperature. This $T_c$ can be identified from the correlation function and from visual inspection of self-similarity. Manipulating the model through a block spin transformation or applying an external field can alter behavior above and below $T_c$, but the critical temperature is ultimately preserved.